\section{Independent audit of the operand-reader reassociation}
\label{sec:orh-independent}
This audit concerns the literal table called \texttt{PA194ORH}. Its dense
transition table has SHA-256
\begin{quote}\small\ttfamily
c16bef7a1587e31adaae747d5bad26721f30bb5c6c33dd4b798be95f5141f1d4.
\end{quote}
The parent is the pinned \texttt{PA194AX11H} table. Native state numbers
are used below; they are names, rather than a count of working states.
The new dense table has 194 working states and 388 binary transition rows.

\subsection{The hidden selector is preserved exactly}
The apparent closed operand routine at state 213 has a second use:
the transition from state 69 on zero moves left into it. Consequently
an arbitrary reassociation of the closed routine cannot be substituted
there without also preserving this history-sensitive entry.

Let the head initially be at an arbitrary position $p$, the state be 69,
and the scanned symbol be zero. Write $h\in\{0,1\}$ for the symbol at
$p-1$. Every transition in the following traces preserves the scanned
symbol. States and head positions are shown after each transition.
For $h=0$ the old trace is
\[
 (69,p)\longrightarrow(213,p-1)\longrightarrow(92,p)
 \longrightarrow(104,p+1),
\]
whereas the new trace is
\[
 (69,p)\longrightarrow(92,p-1)\longrightarrow(111,p)
 \longrightarrow(104,p+1).
\]
For $h=1$ the old and new traces are respectively
\[
 (69,p)\longrightarrow(213,p-1)\longrightarrow(90,p),
 \qquad
 (69,p)\longrightarrow(92,p-1)\longrightarrow(90,p).
\]
In the zero-history case the last transition is legitimate because the
initially scanned cell at $p$ is still zero. Both cases therefore give
exact equality of exit state, head position, and entire tape, for every
initial tape satisfying the single entry condition. The equality does
not depend on a finite register bound or on a blank-tape simulation.
State 111 is an existing reader, and its zero transition is unchanged.

\subsection{The closed operand dispatch is recoded}
The source routine changes from
\[
 ((\mathrm{node32},(\mathrm{node53},\mathrm{node55})),
   \mathrm{node65})
 \quad\hbox{to}\quad
 (\mathrm{node32},(\mathrm{node53},
   (\mathrm{node55},\mathrm{node65}))).
\]
Starting at state 213, the old and new preserved-symbol branch words
reach the same macro entry states as follows.
\begin{center}
\begin{tabular}{ccll}
Old word & New word & Exit state & Source component\\\hline
00 & 0 & 104 & \texttt{node32}\\
010 & 10 & 107 & \texttt{node53}\\
011 & 110 & 98 & \texttt{node55}\\
1 & 111 & 90 & \texttt{node65}
\end{tabular}
\end{center}
The new reader 333 implements the last two branches and is shared with
the AX grammar's existing \texttt{node55};\texttt{node65} pattern.
The finite trace check establishes these four dispatch identities.
Since branch depths can affect register-program jumps, preservation of
the leaf order alone would not suffice. An independent source decoder
also compared the fully resolved instruction lists, as described next.

\subsection{Independent source and unchanged register boxes}
The checker \texttt{ORH\_audit.py} imports neither the construction
compiler nor its symbolic executor. It expands the old and new source
trees and resolves jumps by lexicographic leaf intervals: leaving a
source subtree goes to the first leaf after its interval; continuing
at an ancestor goes to the first leaf in its interval. Decrement's zero
branch uses the former operation at the prescribed zero ancestor.
This independent calculation finds 3,671 sites in each source, with
identical resolved instructions and jump destinations. Exactly 134
source addresses change, and the maximum new address length is 37.
It also compares its output to the construction's stored sites and
instructions.

Below native state number 1300, the only controls with changed rows,
including the new control, are
\[
 \{69,92,185,213,333\}.
\]
The accelerated register-box interface uses the 23 controls
\[
\begin{split}
\{&0,1,2,3,16,22,25,32,33,36,40,48,51,54,55,60,\\
  &64,66,74,94,97,113,198\}.
\end{split}
\]
All 46 rows at these controls are literally identical in the parent and
new tables. Other changed controls are AX readers numbered at least
2000 and likewise lie outside that interface. Thus the change does not
silently change a primitive register-box contract used by the inherited
symbolic simulation.

\subsection{Scope of the conclusion}
The independent receipt is \texttt{ORH\_audit.json}. It certifies the
finite selector traces, the four dispatch cases, the unchanged box rows,
and equality of the complete resolved source programs. Separately, the
construction receipt reports all 4,778 symbolic source contracts and the
entire blank initialization frame; those computations are not rerun by
this independent checker. The reported maximum 506 belongs to the
symbolic executor, which accelerates tape runs and invokes register-box
contracts. It is not a uniform 506-step bound on concrete Turing-machine
execution for arbitrary natural registers.

The appropriate correctness transfer identifies source instruction
ordinals and translates their physical program addresses. It is not an
assertion that old and new complete machines take equal numbers of
physical steps. Once the inherited source-level calibration and the
new physical contracts are used, this gives the same refutation search
as the parent. A finite Python audit is not itself a new Lean proof of
the absolute iPRA calibration.
